\documentclass[11pt]{article}

\usepackage[a4paper,margin=28mm]{geometry}
\usepackage{amsmath,amssymb,amsthm,mathtools}
\usepackage{microtype}
\usepackage[hidelinks]{hyperref}

\newtheorem{theorem}{Theorem}
\newtheorem{proposition}[theorem]{Proposition}
\newtheorem{corollary}[theorem]{Corollary}
\theoremstyle{remark}
\newtheorem{remark}[theorem]{Remark}

\newcommand{\Z}{\mathbb{Z}}
\newcommand{\Q}{\mathbb{Q}}

\title{A Polynomial Parametrization of the Surface\\[2mm]
\texorpdfstring{$x^3+yz+1=0$}{x^3+yz+1=0}}
\author{}
\date{}

\begin{document}
\maketitle
\vspace{-1.5em}

\begin{abstract}
We construct an explicit two-parameter polynomial family of integral
solutions to
\[
 x^3+yz+1=0.
\]
More precisely, the polynomial map
\[
 (a,b)\longmapsto
 \bigl(ab-1,\ a,\ -b(a^2b^2-3ab+3)\bigr)
\]
takes $\Z^2$ into the set of integral solutions.  Over $\Q$, the same
map is a bijection onto the full set of rational solutions.  We prove
these assertions, characterize exactly which integral points arise
from integral parameter values, record two useful one-parameter
families, and explain in detail the factorization and
parameter-splitting idea that leads naturally to the construction.
\end{abstract}

\section{The polynomial construction}

For a commutative ring $R$, write
\[
 S(R)=\bigl\{(x,y,z)\in R^3:x^3+yz+1=0\bigr\}.
\]
The elementary identity underlying the construction is
\begin{equation}\label{eq:factorization}
 X^3+1=(X+1)(X^2-X+1).
\end{equation}

\begin{theorem}[Two-parameter polynomial family]\label{thm:main}
For arbitrary $a,b\in\Z$, define
\begin{equation}\label{eq:parametrization}
 \boxed{
 \begin{aligned}
 x&=ab-1,\\
 y&=a,\\
 z&=-b\bigl(a^2b^2-3ab+3\bigr).
 \end{aligned}}
\end{equation}
Then $(x,y,z)\in S(\Z)$.  Equivalently, the following identity holds
in the polynomial ring $\Z[a,b]$:
\[
 (ab-1)^3
 +a\bigl[-b(a^2b^2-3ab+3)\bigr]+1=0.
\]
\end{theorem}

\begin{proof}
Set $X=ab-1$.  Then
\[
 X+1=ab
\]
and
\begin{align*}
 X^2-X+1
 &=(ab-1)^2-(ab-1)+1\\
 &=a^2b^2-3ab+3.
\end{align*}
Applying \eqref{eq:factorization} gives
\[
 X^3+1
 =ab\bigl(a^2b^2-3ab+3\bigr).
\]
On the other hand, the definitions of $y$ and $z$ give
\[
 yz
 =a\left[-b\bigl(a^2b^2-3ab+3\bigr)\right]
 =-ab\bigl(a^2b^2-3ab+3\bigr).
\]
Hence $X^3+yz+1=0$.  Since all three coordinates in
\eqref{eq:parametrization} are polynomials with integer coefficients,
they are integers whenever $a$ and $b$ are integers.
\end{proof}

\begin{corollary}[Infinitely many integral solutions]\label{cor:infinite}
For every $t\in\Z$,
\begin{equation}\label{eq:oneparameter}
 \boxed{(x,y,z)=\bigl(t,\ t+1,\ -(t^2-t+1)\bigr)}
\end{equation}
is an integral solution.  These solutions are pairwise distinct as
$t$ varies; in particular, $S(\Z)$ is infinite.
\end{corollary}

\begin{proof}
In Theorem~\ref{thm:main}, take $a=t+1$ and $b=1$.  Then
$x=(t+1)-1=t$, while
\[
 a^2-3a+3=(t+1)^2-3(t+1)+3=t^2-t+1,
\]
which gives \eqref{eq:oneparameter}.  Distinct values of $t$ give
distinct values of the first coordinate $x=t$.
\end{proof}

\section{Completeness over the rational numbers}

The polynomial family above is not merely a convenient source of
solutions.  With rational parameters it describes every rational
point on the surface.

\begin{theorem}[All rational solutions]\label{thm:rational}
Define
\[
 \Phi:\Q^2\longrightarrow S(\Q),\qquad
 \Phi(a,b)=
 \bigl(ab-1,\ a,\ -b(a^2b^2-3ab+3)\bigr).
\]
Then $\Phi$ is a bijection.  Its inverse is given explicitly by
\begin{equation}\label{eq:inverse}
 \Phi^{-1}(x,y,z)=
 \begin{cases}
 \displaystyle\left(y,\frac{x+1}{y}\right),&y\ne0,\\[3mm]
 \displaystyle\left(0,-\frac{z}{3}\right),&y=0.
 \end{cases}
\end{equation}
\end{theorem}

\begin{proof}
Theorem~\ref{thm:main} proves the defining relation as a polynomial
identity, so it remains valid for $a,b\in\Q$.  Thus $\Phi$ is
well-defined.

Let $(x,y,z)\in S(\Q)$.  First suppose that $y\ne0$, and put
\[
 a=y,
 \qquad
 b=\frac{x+1}{y}.
\]
Then $ab=x+1$, and hence $ab-1=x$.  Moreover,
\begin{align*}
 a^2b^2-3ab+3
 &=(x+1)^2-3(x+1)+3\\
 &=x^2-x+1.
\end{align*}
Consequently,
\begin{align*}
 -b(a^2b^2-3ab+3)
 &=-\frac{x+1}{y}(x^2-x+1)\\
 &=-\frac{x^3+1}{y}\\
 &=z,
\end{align*}
where the last equality is exactly the equation
$x^3+yz+1=0$.  Hence $\Phi(a,b)=(x,y,z)$.

Now suppose that $y=0$.  The defining equation becomes $x^3+1=0$.
Over $\Q$ this forces $x=-1$.  Taking
\[
 a=0,
 \qquad
 b=-\frac{z}{3},
\]
we obtain
\[
 \Phi(a,b)=(-1,0,-3b)=(-1,0,z).
\]
Thus $\Phi$ is surjective and \eqref{eq:inverse} is a right inverse.

To prove injectivity, suppose $\Phi(a,b)=\Phi(a',b')$.  Equality of the
second coordinates gives $a=a'$.  If this common value is nonzero,
equality of the first coordinates gives $ab=a'b'$, hence $b=b'$.
If $a=a'=0$, equality of the third coordinates gives
$-3b=-3b'$, hence again $b=b'$.  Therefore $(a,b)=(a',b')$, and
$\Phi$ is bijective.
\end{proof}

\begin{proposition}[The image of integral parameter values]\label{prop:integralimage}
Let $(x,y,z)\in S(\Z)$.  Then $(x,y,z)\in\Phi(\Z^2)$ if and only if
one of the following mutually exclusive conditions holds:
\begin{enumerate}
 \item $y\ne0$ and $y\mid x+1$;
 \item $y=0$ and $3\mid z$.
\end{enumerate}
In the first case the unique integral parameters are
\[
 (a,b)=\left(y,\frac{x+1}{y}\right),
\]
and in the second case they are
\[
 (a,b)=\left(0,-\frac{z}{3}\right).
\]
\end{proposition}

\begin{proof}
By Theorem~\ref{thm:rational}, every point of $S(\Z)$ has the unique
rational parameters displayed in \eqref{eq:inverse}.  If $y\ne0$,
those parameters are integral exactly when
$(x+1)/y\in\Z$, which is equivalent to $y\mid x+1$.  If $y=0$, the
equation forces $x=-1$, and the unique parameters are integral exactly
when $-z/3\in\Z$, which is equivalent to $3\mid z$.
\end{proof}

\begin{remark}\label{rem:notallintegral}
The distinction in Proposition~\ref{prop:integralimage} is essential.
For example, $(2,9,-1)$ is an integral solution because
$2^3+9(-1)+1=0$, but its unique parameters are $(9,1/3)$, not an
integer pair.  Thus \eqref{eq:parametrization} gives a large polynomial
family of integral points but does not exhaust $S(\Z)$ when the
parameters are required to be integral.
\end{remark}

\section{How the construction was found}\label{sec:discovery}

The construction is guided by two requirements: the equation should
hold identically, and every displayed coordinate should remain a
polynomial in freely chosen parameters.  The following sequence of
observations leads directly to \eqref{eq:parametrization}.

\subsection*{1. Isolate the product}
The equation may be rewritten as
\begin{equation}\label{eq:product}
 yz=-(x^3+1).
\end{equation}
Thus one does not need to control $y$ and $z$ separately at first; one
only needs to split the polynomial $-(x^3+1)$ into two factors.

\subsection*{2. Use the sum-of-cubes factorization}
The cubic has the canonical factorization
\[
 x^3+1=(x+1)(x^2-x+1).
\]
The most immediate choice in \eqref{eq:product} is therefore
\[
 y=x+1,
 \qquad
 z=-(x^2-x+1).
\]
This already gives the one-parameter family
\[
 (x,y,z)=\bigl(t,t+1,-(t^2-t+1)\bigr).
\]
It is correct, but it uses only one free parameter.

\subsection*{3. Split the linear factor to create a second parameter}
To introduce a second independent parameter without creating
fractions, the linear factor $x+1$ is the natural place to intervene.
Write
\begin{equation}\label{eq:split}
 x+1=ab.
\end{equation}
This equation can be solved polynomially for $x$:
\[
 x=ab-1.
\]
We then distribute the two factors $a$ and $b$ between $y$ and $z$:
\[
 y=a,
 \qquad
 z=-b(x^2-x+1).
\]
With these choices,
\[
 yz=-ab(x^2-x+1)=-(x+1)(x^2-x+1)=-(x^3+1),
\]
so the original equation is forced.

Finally, substituting $x=ab-1$ into the remaining quadratic factor
gives
\[
 x^2-x+1
 =(ab-1)^2-(ab-1)+1
 =a^2b^2-3ab+3.
\]
This is precisely the polynomial appearing in
\eqref{eq:parametrization}.

\subsection*{4. Why this split is particularly efficient}
One could try to split the quadratic factor instead, for example by
imposing $x^2-x+1=ab$.  However, solving that relation for $x$ would
normally require extracting a square root, since
\[
 (2x-1)^2=4ab-3.
\]
That approach therefore introduces an additional square condition and
does not immediately produce polynomial coordinates.  By contrast,
$x+1=ab$ is linear in $x$, so eliminating $x$ preserves polynomiality
and costs no auxiliary condition.  The useful design principle is:
\emph{split a factor that is linear in the original variable, then
absorb the complementary factor into the remaining product variable.}

\subsection*{5. Why rational completeness is visible from the same idea}
Suppose a rational solution has $y\ne0$.  Since the construction sets
$a=y$ and requires $ab=x+1$, the only possible value of the second
parameter is
\[
 b=\frac{x+1}{y}.
\]
The original equation then forces the $z$-coordinate to equal the one
produced by the parametrization.  Thus the factor-splitting argument
not only creates solutions; it also recovers every rational solution
away from $y=0$.  The remaining locus $y=0$ is the line
$(-1,0,z)$, and the specialization $a=0$ gives
$\Phi(0,b)=(-1,0,-3b)$, which covers that line over $\Q$ as well.

\section{A second one-parameter family}

There is also a symmetric-looking family obtained from a different
factorization pattern.

\begin{proposition}\label{prop:symmetric}
For every $t\in\Z$,
\begin{equation}\label{eq:symmetric}
 \boxed{(x,y,z)=\bigl(-t^2,\ t^3-1,\ t^3+1\bigr)}
\end{equation}
is an integral solution of $x^3+yz+1=0$.
\end{proposition}

\begin{proof}
Using the difference-of-squares factorization,
\[
 yz=(t^3-1)(t^3+1)=t^6-1.
\]
Since $x=-t^2$, one has $x^3=-t^6$.  Therefore
\[
 x^3+yz+1=-t^6+(t^6-1)+1=0.
\]
\end{proof}

\begin{remark}[Intuition for the second family]
To find \eqref{eq:symmetric}, choose $x=-t^2$ so that $-x^3=t^6$ is a
perfect square and
\[
 -(x^3+1)=t^6-1.
\]
The right-hand side factors as $(t^3-1)(t^3+1)$, which can be assigned
directly to $y$ and $z$.  This illustrates the same general strategy:
make $-(x^3+1)$ factor in a parameter-controlled way and assign its
factors to the product $yz$.
\end{remark}

\section{Conclusion}

The principal polynomial parametrization is
\[
 \boxed{
 (x,y,z)=
 \bigl(ab-1,\ a,\ -b(a^2b^2-3ab+3)\bigr),
 \qquad a,b\in\Z.}
\]
Its validity is an identity in $\Z[a,b]$, so no congruence,
divisibility, or exceptional-case condition is required for integral
parameter values.  Over $\Q$, the same map is a bijection onto the
entire rational solution set, with inverse given by
\eqref{eq:inverse}.  The construction follows naturally from
factoring $x^3+1$ and splitting its linear factor $x+1$ into two free
parameters.

\end{document}
